% State of the Art section. Paste into your own thesis, or \input it.
%
% Built from docs_thesis/find_research.md \S6 "The gap, stated for \S4" and
% docs_thesis/literature.md Axis 4's recommended framing for objective 9.
% Verification status: zamfir2024anomalies and rudnicka2022quartz DOI-verified against Crossref
% this session. The rest ([BIBLIOGRAPHY VERIFIED]) per docs_thesis/find_research.md and
% docs_thesis/literature.md.

\section{The gap}
\label{sec:the-gap}

The pieces this chapter has assembled point at one omission rather than at open ground everywhere.
The mechanism from shape to plaque is settled at lesion scale, and the shear field it runs through is
recoverable from geometry alone with its prognostic content intact
(Section~\ref{sec:shape-enough}).
Individual tree-scale features --- bifurcation angle, dominance, tortuosity, branch presence --- each
carry a small evidence base pointing in a consistent direction (Section~\ref{sec:geometry}).
The outcome cohorts that could test those features at population scale already exist and are already
published, including the one this project moves to: a population cohort with no known heart
disease, screened by CCTA, carrying adjudicated myocardial infarction \cite{fuchs2023cgps}.
What does not exist is the join.
No study measures tree-scale geometry across a population of thousands, describes its distribution,
and asks which individuals fall outside it (Section~\ref{sec:scale-gap}).

That is not for lack of machine-learning effort in the field, only effort spent elsewhere.
It went to plaque radiomics, where ten studies pooled to an AUC near 0.80 predicting major adverse
cardiac events from lesion texture and morphology \cite{ma2025mlmace}, and its graph neural networks
model patient similarity or predict centerline segment labels rather than the tree's own topology
(Section~\ref{sec:nearest-competitors}) \cite{yaseliani2025gnn}.

The clinical literature is not empty-handed on the word "abnormal," but what it holds is a different
operation from the one objective 9 asks for.
A coronary anomaly, in cardiology, is a morphological feature of a vessel's origin, course or
termination, defined categorically and diagnosed the way a radiologist recognizes a named pattern:
detected in 192 of 8733 CCTA scans (2.2\%) at one center before 2020 and 158 of 4025 (3.9\%)
afterward \cite{zamfir2024anomalies}.
That is a taxonomy, not a distribution --- the absent left main variant this thesis already treats as
a named anatomical category (Section~\ref{sec:variants}) is exactly this kind of entry.
It is not the operation of treating branching pattern, bifurcation angle and tortuosity as a joint
distribution over an unselected population and flagging which values sit in its tails, regardless of
whether any of them has a name.
Framed this way, objective 9 is the unsupervised, data-driven counterpart to a concept clinical
practice already has, not a claim to unoccupied territory --- a narrower and more defensible position
than a bare novelty claim, and one that supplies its own validation handle.
A method built to flag distributional outliers should rediscover the known categories in this cohort
without being told about them: the 11 absent-left-main sides and the 41 left-dominant, 30 codominant
hearts already recorded in ImageCAS-X \cite{bransby2026imagecasx}.
A method that misses those is not trustworthy on the cases where the answer is not already known
either.

The nearest methodological precedent sits in a different organ.
Rudnicka et al.\ trained the QUARTZ algorithm on retinal vessel width, area and tortuosity in 88{,}052
UK Biobank participants and used it, combined with clinical risk factors, to predict circulatory
death, myocardial infarction and stroke; the retinal vasculature alone captured half to two-thirds of
circulatory deaths among those it flagged highest-risk \cite{rudnicka2022quartz}.
That is population-scale vascular topology used exactly the way objective 9 proposes to use the
coronary tree, in a vascular bed that is easier to image but carries none of the direct clinical
stakes.
No coronary-specific equivalent was found in any of the searches this chapter is built on
(\texttt{docs\_thesis/find\_research.md}, \texttt{docs\_thesis/literature.md}).
That is the gap this thesis occupies.
